\chapter{Úvod}

\section{Motivace}
Téma bakalářské práce bylo zvoleno na~základě potřeb studentského raketového spolku Czech Rocket Society jak z~hlediska popularizace kosmonautiky v~České republice, tak pro~vnitřní fungování spolku.
Platforma je navržena jako nástroj, který umožní lépe prezentovat reálné schopnosti studentů působících v~kosmickém raketovém sektoru, včetně možnosti zobrazovat živá telemetrická data z~letů raket. Součástí jsou také menší podpůrné aplikace pro každodenní činnost spolku.
Cílem je vytvořit modulární platformu složenou z~několika subsystémů, z~nichž každý podporuje konkrétní spolkovou činnost a~činí ji efektivnější.
Aplikace přináší nejen nové funkcionality, ale také částečně řeší problémy stávajícího webu. Ten je postaven na~technologiích, které neumožňují jednoduchou údržbu obsahu ani jeho aktualizace bez hlubší znalosti programovacích či značkovacích jazyků a~základů systémů správy verzí.
Po~designové stránce by měl web splňovat nároky spolku, obstát ve~srovnání s~konkurencí a~vhodně propagovat kosmonautiku v~České republice.

\section{Struktura práce}
Práce je rozdělena do následujících kapitol:
\begin{itemize}
    \item \textbf{Cíl a~metodika práce} -- vymezení cílů, popis zvolené metodiky a~využití nástrojů umělé inteligence.
    \item \textbf{Teoretická část a~analýza} -- analýza domény, přehled existujících řešení a~teoretický základ použitých technologií.
    \item \textbf{Návrh řešení} -- architektura systému, datový model, návrh API (Application Programming Interface) a~uživatelského rozhraní.
    \item \textbf{Implementace} -- popis realizace jednotlivých částí platformy.
    \item \textbf{Testování} -- metodika testování a~shrnutí výsledků.
    \item \textbf{Shrnutí a~diskuse výsledků} -- porovnání s~výchozím stavem, soulad výsledků s~literaturou a~okolnosti ovlivňující průběh práce.
    \item \textbf{Závěry a~doporučení} -- zhodnocení splnění cílů a~možnosti budoucího rozvoje.
\end{itemize}